\documentclass[9pt,a4paper]{extarticle}
\usepackage{parskip}

\usepackage[utf8]{inputenc}
\usepackage[T1]{fontenc}
\usepackage[brazil]{babel}
\usepackage{amsmath, amssymb, amsthm}
\usepackage{geometry}
\usepackage{xcolor}
\usepackage{hyperref}
\usepackage{booktabs}
\usepackage{enumitem}
\usepackage{microtype}
\usepackage{csquotes}
\usepackage{mdframed}
\usepackage{pgffor}
\usepackage{tabto}

\newcommand{\answerlines}[1]{%
  \par\vspace{4pt}%
  \foreach \n in {1,...,#1}{%
    \noindent\makebox[\linewidth]{\dotfill}\\[6pt]%
  }%
}

\definecolor{important}{RGB}{255, 100, 100}
\definecolor{notice}{RGB}{0, 102, 204}
\definecolor{headerblue}{RGB}{152, 38, 73}
\definecolor{accentblue}{RGB}{113, 162, 182}

\newcommand{\impt}[1]{\textcolor{important}{\textbf{#1}}}
\newcommand{\ntc}[1]{\textcolor{notice}{#1}}

\setlist{nosep}
\geometry{margin=2.5cm, top=1.5cm}

\newmdenv[
  linecolor=headerblue,
  backgroundcolor=headerblue!6,
  linewidth=1.2pt,
  innerleftmargin=10pt,
  innerrightmargin=10pt,
  innertopmargin=8pt,
  innerbottommargin=8pt,
  skipabove=6pt,
  skipbelow=6pt
]{notebox}

\newmdenv[
  linecolor=accentblue,
  backgroundcolor=accentblue!8,
  linewidth=0.8pt,
  innerleftmargin=10pt,
  innerrightmargin=10pt,
  innertopmargin=8pt,
  innerbottommargin=8pt,
  skipabove=6pt,
  skipbelow=6pt
]{casobox}

\title{%
  \vspace{-1cm}%
  \textcolor{headerblue}{\textbf{Prova 1}}\\[4pt]
  \large Metodologia Científica para Ciência da Computação --- BCC502/DECOM/UFOP \\
  }
\date{}

\begin{document}
\maketitle
\vspace{-1.5cm}
Nome:................................................................................................... Matrícula: ...................................

% ==================================================================
1. Considere o seguinte cenário, próximo do que vocês podem encontrar na prática:
 
\begin{casobox}
Um grupo de estudantes desenvolve um algoritmo que afirmam ser mais rápido que os algoritmos clássicos para um certo problema de ordenação. Publicam o resultado com testes em algumas entradas. Um colega aponta que as entradas escolhidas favorecem o método. Os autores publicam uma versão revisada: introduzem um pré-processamento ``essencial'' que melhora o tempo nos casos contestados. Outro grupo tenta reproduzir os resultados em sua própria máquina e obtém tempos piores que os clássicos; os autores respondem que faltou ajustar um parâmetro de cache específico do hardware deles. Um terceiro grupo testa o algoritmo em entradas novas, geradas aleatoriamente; o algoritmo é mais lento; os autores respondem que essas entradas têm uma ``estrutura atípica'' e que o método funciona bem para entradas ``realistas''. A cada nova dificuldade, o grupo original introduz uma modificação que protege a alegação original.
\end{casobox}
 
\ntc{Lembrete:} ao responder os itens abaixo, espera-se uso preciso de termos como \emph{cinturão protetor}, \emph{núcleo firme}, modificação \emph{ad hoc}, mudança progressiva/degenerativa, predição nova corroborada, predição arriscada, heurística posisitiva / negativa  --- onde se aplicarem.
 
\vspace{0.2cm}
 
\noindent\textbf{(a)} Sob a ótica de Popper (falsificacionismo sofisticado) e de Lakatos, há um \impt{vocabulário técnico} bastante específico (veja lembrete) para descrever o que está acontecendo. Use esse vocabulário e diga \impt{o que está dando errado} aqui --- sem usar a palavra ``ruim'' ou equivalentes vagos.
 
\answerlines{26}
 
\noindent\textbf{(b)} Um defensor do grupo pode responder: \emph{``vocês estão sendo falsificacionistas ingênuos --- a primeira refutação não derruba um programa''}. Essa defesa é \impt{legítima} à luz de Lakatos? Sob quais condições ela passaria a ser \impt{ilegítima}? Sua resposta deve identificar o critério lakatosiano envolvido, não apenas dizer ``depende''.
 
\answerlines{18}
 
\vspace{-0.5cm}
\noindent\textbf{(c)} Agora mude de lente: descreva o mesmo episódio do ponto de vista de Kuhn. O que um kuhniano \impt{veria de diferente} que um lakatosiano não vê? Você pode usar conceitos como \emph{ciência normal}, \emph{quebra-cabeças}, \emph{exemplares}, \emph{valores compartilhados} --- mas só os que de fato se aplicam ao caso.
 
\answerlines{18}

% ==================================================================
2. Cada uma das afirmações abaixo é uma \impt{distorção plausível} de uma ideia que estudamos --- o tipo de coisa que ``parece certo'' mas confunde dois conceitos próximos. Para cada uma, identifique a confusão e formule a versão correta. Não se prenda à definição do Chalmers; explique \impt{por que a distinção importa}.

\begin{enumerate}[label=(\alph*), leftmargin=*]
  \item \emph{``Como paradigmas rivais são incomensuráveis, não há como cientistas de paradigmas diferentes discutirem entre si --- eles literalmente não se entendem.''}
  \answerlines{19}

  \vspace{-0.3cm}
  \item \emph{``O núcleo firme de um programa de pesquisa lakatosiano é, na prática, irrefutável --- ou seja, é tratado pelos cientistas como verdade estabelecida.''}
  \answerlines{19}

  \item \emph{``A diferença entre uma modificação legítima e uma modificação \emph{ad hoc} é que a primeira é introduzida por uma boa razão e a segunda apenas para salvar a teoria.''}
  \answerlines{19}
 \vspace{-0.3cm}
  \item \emph{``Como toda observação é carregada de teoria, qualquer experimento pode ser interpretado de qualquer jeito --- a falsificação é, no fim, uma escolha sociológica.''}
  \answerlines{19}
\end{enumerate}

\end{document}
